\section{Final Exam: From Beginner to Ready}
\chapterline

This exam is not a barrier erected to keep you out of higher mathematics. It is a diagnostic mirror. 

When you solve an engineering problem, you check your answer against numerical intuition or physical units. But when you write a pure mathematical proof, your only validator is the relentless internal consistency of your logical chain. This final exam is designed to test whether your logical instincts have become automatic: whether you instinctively verify types, untangle alternating quantifiers, distinguish discovery from deduction, and respect the boundaries of definitions.

\begin{historybox}
In August 1900, at the International Congress of Mathematicians in Paris, David Hilbert delivered a historic address that shaped the twentieth century. He presented a list of 23 unsolved problems, challenging mathematicians to settle the foundations of mathematics, arithmetic, and geometry. Hilbert believed passionately in the ultimate consistency and completeness of human reason, proclaiming his famous conviction: \emph{``Wir m\"ussen wissen, wir werden wissen''} (``We must know, we will know'').

The quest to build an unshakeable foundation for mathematics brought together the finest minds of the modern era---Gottlob Frege, Bertrand Russell, Henri Poincar\'e, Ernst Zermelo, John von Neumann, and Kurt G\"odel. The grammar of proof you have learned in this book is the direct fruit of their struggle. By uniting logic, sets, functions, quotients, and induction, they turned mathematics from a collection of calculation recipes into a unified deductive universe.
\end{historybox}

\subsection{The Constellation of Mathematical Readiness}

Before you begin, look at the interconnected web of ideas you have traversed. None of these topics exist in isolation; they form a single coherent architecture:

\begin{center}
\begin{tikzpicture}[scale=1.1, >=Stealth]
  \tikzset{
    corenode/.style={draw=structurepurple!90!black, fill=structurepurple!12, very thick, rounded corners=4pt, inner sep=6pt, font=\normalsize\bfseries, align=center},
    pillarnode/.style={very thick, rounded corners=3pt, inner sep=5pt, font=\small\bfseries, align=center},
    fieldnode/.style={draw=gray!60, fill=gray!8, thick, rounded corners=3pt, inner sep=4pt, font=\footnotesize, align=center},
    conedge/.style={draw=gray!50, thick},
    radedge/.style={->, draw=structurepurple!80!black, thick, shorten >=2pt, shorten <=2pt}
  }

  % Central Hub
  \node[corenode] (center) at (5.0, 3.2) {Proof-Based\\Mathematical\\Maturity};

  % Surrounding Pillars
  % Top Left: Logic
  \node[pillarnode, draw=sourceblue!80!black, fill=sourceblue!15] (logic) at (1.5, 5.4) {1. Logic \&\\Quantifiers};
  % Top Right: Sets
  \node[pillarnode, draw=targetgreen!80!black, fill=targetgreen!15] (sets) at (8.5, 5.4) {2. Sets \&\\Boolean Duality};
  % Bottom Left: Induction
  \node[pillarnode, draw=trapred!80!black, fill=trapred!15] (induct) at (1.5, 1.0) {4. Induction \&\\Well-Ordering};
  % Bottom Right: Quotients
  \node[pillarnode, draw=quotientorange!90!black, fill=quotientorange!18] (quot) at (8.5, 1.0) {3. Functions,\\Relations, Quotients};

  % Edges between pillars and center
  \draw[conedge, draw=sourceblue!70] (logic) -- (center);
  \draw[conedge, draw=targetgreen!70] (sets) -- (center);
  \draw[conedge, draw=trapred!70] (induct) -- (center);
  \draw[conedge, draw=quotientorange!70] (quot) -- (center);

  % Outer Ring: Advanced Subjects
  \node[fieldnode] (anal) at (5.0, 6.0) {Real Analysis\\$(\eps\text{--}\delta \text{ Contracts})$};
  \node[fieldnode] (alg) at (10.2, 3.2) {Abstract Algebra\\$(G/\ker\phi \cong \operatorname{im}\phi)$};
  \node[fieldnode] (top) at (5.0, 0.4) {Topology\\$(f^{-1}(V) \text{ Open})$};
  \node[fieldnode] (geom) at (-0.2, 3.2) {Differential \& Alg.\\Geometry};

  % Radiating arrows from center to advanced fields
  \draw[radedge] (center) -- (anal);
  \draw[radedge] (center) -- (alg);
  \draw[radedge] (center) -- (top);
  \draw[radedge] (center) -- (geom);
\end{tikzpicture}
\end{center}

\subsection{Controversies, False Starts, and the Foundational Crisis}

The road to modern mathematical rigor was forged through an intense existential crisis (*Grundlagenkrise*):

\begin{itemize}
  \item \textbf{Frege's Heartbreak and Russell's Postcard (1902)}: Gottlob Frege dedicated decades of his life to constructing arithmetic from pure logic in his two-volume opus, \emph{Grundgesetze der Arithmetik}. In 1902, as the second volume was in press, he received a short letter from the young Bertrand Russell pointing out that unrestricted set comprehension allows the set $R = \{x : x \notin x\}$, leading directly to $R \in R \iff R \notin R$. Frege added a famous, poignant appendix: ``Hardly anything more unwelcome can befall a scientific writer than to have one of the foundations of his edifice shaken after the work is finished.'' Logic had to be rebuilt with strict type and set boundaries.
  \item \textbf{G\"odel's Incompleteness Shock (1931)}: Hilbert believed that every mathematical statement could either be proven or disproven by a consistent, mechanical formal system. In 1931, a 25-year-old Austrian logician named Kurt G\"odel shattered Hilbert's program. G\"odel's First Incompleteness Theorem proved that \emph{any} consistent axiomatic system rich enough to express arithmetic contains true statements that cannot be proven within the system. His Second Incompleteness Theorem proved that such a system cannot even prove its own consistency. Pure mathematics is not a closed mechanical algorithm; it is an open-ended ocean of truth.
  \item \textbf{The Syntax vs.\ Semantics Confusion}: A persistent false start among aspiring mathematicians is assuming that manipulating symbols correctly is the same as understanding mathematical truth. Formal syntax is merely the recording medium; true mathematical maturity lies in understanding the semantic invariants and structures behind the symbols.
\end{itemize}

\subsection{Part 1: The Grammar of Statements and Quantifiers}

\begin{enumerate}
  \item Explain the grammatical difference between a \emph{statement} and a \emph{predicate}. Give an explicit example of a predicate $P(x, y)$ that becomes a true statement under one order of quantifiers ($\forall x \exists y$) but a false statement when the quantifiers are reversed ($\exists y \forall x$).
  \item Negate each of the following statements completely, simplifying so that no negation symbol appears outside a quantifier or connective:
  \begin{enumerate}
    \item $\forall \eps > 0, \; \exists \delta > 0, \; \forall x \in \R, \quad (0 < |x - c| < \delta \implies |f(x) - L| < \eps)$.
    \item For every natural number $n \ge 2$, if $n$ is not prime, then $n$ possesses a prime divisor $p \le \sqrt{n}$.
  \end{enumerate}
  \item State the contrapositive and the converse of the implication: ``If a sequence $(a_n)$ is unbounded, then it does not converge.'' Which of these two is logically equivalent to the original statement?
\end{enumerate}

\subsection{Part 2: Sets, Inclusions, and Type Integrity}

\begin{enumerate}[resume]
  \item Prove that for any three sets $A, B$, and $C$:
  \[
  A \setminus (B \cup C) = (A \setminus B) \cap (A \setminus C).
  \]
  Structure your proof by demonstrating that an arbitrary element $x$ belongs to the left-hand side if and only if it belongs to the right-hand side.
  \item Let $A$ and $B$ be arbitrary sets. Prove or disprove each claim:
  \begin{enumerate}
    \item $\pow(A \cap B) = \pow(A) \cap \pow(B)$.
    \item $\pow(A \cup B) = \pow(A) \cup \pow(B)$.
  \end{enumerate}
  If a claim is false, provide an explicit, finite counterexample.
\end{enumerate}

\subsection{Part 3: Functions, Morphisms, and Quotients}

\begin{enumerate}[resume]
  \item Let $f: X \to Y$ be a function, and let $S, T \subseteq Y$.
  \begin{enumerate}
    \item Prove that $f^{-1}(S \cap T) = f^{-1}(S) \cap f^{-1}(T)$.
    \item Provide a simple counterexample showing that forward images do not preserve intersections: find a function $g$ and subsets $A, B$ such that $g(A \cap B) \neq g(A) \cap g(B)$. What condition on $g$ restores equality?
  \end{enumerate}
  \item Define a relation on $\Z \times (\Z \setminus \{0\})$ by $(a, b) \sim (c, d) \iff ad = bc$.
  \begin{enumerate}
    \item Prove that $\sim$ is an equivalence relation.
    \item Explain how the quotient set $(\Z \times (\Z \setminus \{0\}))/\sim$ constructs the rational numbers $\Q$.
  \end{enumerate}
  \item Let $\phi: G \to H$ be a group homomorphism. Prove that $\ker(\phi)$ is a normal subgroup of $G$: that is, for every $n \in \ker(\phi)$ and every $g \in G$, show that $gng^{-1} \in \ker(\phi)$.
\end{enumerate}

\subsection{Part 4: Induction and the Minimal Counterexample}

\begin{enumerate}[resume]
  \item Prove by mathematical induction that for all integers $n \ge 1$:
  \[
  1^2 + 2^2 + 3^2 + \dots + n^2 = \frac{n(n+1)(2n+1)}{6}.
  \]
  Clearly state your induction hypothesis and identify exactly where it is invoked.
  \item Use the Well-Ordering Principle (the method of minimal counterexample) to prove that every integer $n > 1$ is divisible by at least one prime number.
\end{enumerate}

\subsection{Diagnostic Self-Assessment Rubric}

To evaluate your proofs, review each response against these four criteria:

\begin{center}
\small
\renewcommand{\arraystretch}{1.25}
\begin{tabular}{p{2.2cm}p{6.5cm}p{6.5cm}}
\hline
\textbf{Check} & \textbf{What to Look For} & \textbf{Common Warning Signs} \\
\hline
\textbf{1. Types} & Are elements, sets, and functions kept distinct? & Writing $x \in f$ or treating $S \subseteq A$ as $S \in A$. \\
\textbf{2. Scope} & Are arbitrary variables introduced before use? & Using $x$ or $\eps$ without saying ``Let $\eps > 0$''. \\
\textbf{3. Order} & Does the proof run strictly forward? & Writing backwards scratchwork as deductive steps. \\
\textbf{4. Definitions} & Are definitions opened and closed properly? & Relying on colloquial words instead of mathematical definitions. \\
\hline
\end{tabular}
\end{center}

\subsection{What Readiness Unlocks---and the Horizon Ahead}

Completing this curriculum unlocks complete readiness for the upper-division landscape of pure mathematics:
\begin{enumerate}
  \item \textbf{Fearless Fluency}: You can open any standard advanced mathematics text without disorientation. When an author writes ``assume for contradiction,'' ``let $\eps > 0$,'' or ``pass to the quotient ring,'' you recognize the exact machinery in motion.
  \item \textbf{Cognitive Armor Against Fallacies}: You are immune to the classic blunders that trap novices: scoping leaks, quantifier transposition, illicit cancellation of zero divisors, and the conflation of subsets with elements.
  \item \textbf{Unified Structural Perspective}: You see that modern mathematics is not a fragmented collection of isolated subjects, but a single conversation. Groups measure symmetries, topologies measure nearness, manifolds measure curvature, and quotients build new worlds.
\end{enumerate}

\paragraph{Honest Downsides and Limitations.}
Yet foundational readiness is only the gateway:
\begin{itemize}
  \item \textbf{Grammar Is Not Literature}: Mastering the grammar of proof gives you fluency, but it does not write the story. Each advanced field (complex analysis, algebraic topology, differential geometry) has its own deep domain intuition, technical lemmas, and geometric visualization that require months of dedicated study.
  \item \textbf{The Incompleteness Horizon}: G\"odel's theorem guarantees that no matter how far pure mathematics advances, there will always be true mathematical statements that transcend any fixed axiomatic framework. Certainty is not a mechanical cage; it is a living quest.
\end{itemize}

\paragraph{Where We Go Next in Pure Mathematics.}
You are now ready to advance into standard upper-division pure mathematics textbooks:
\begin{itemize}
  \item Daniel J. Velleman's \emph{How to Prove It} or Richard Hammack's \emph{Book of Proof} for further practice in formal deduction.
  \item Walter Rudin's \emph{Principles of Mathematical Analysis} or Stephen Abbott's \emph{Understanding Analysis} for rigorous real analysis.
  \item David S. Dummit and Richard M. Foote's \emph{Abstract Algebra} or Michael Artin's \emph{Algebra} for modern structural algebra.
  \item James Munkres' \emph{Topology} for point-set and algebraic topology.
  \item Interactive Theorem Provers (such as Lean 4 and Coq), where proofs are compiled directly into computer-checked mathematical certainty.
\end{itemize}

\begin{intuitionbox}
You have crossed the chasm from calculating answers to giving reasons. The entire architecture of pure mathematics is open before you.
\end{intuitionbox}

\newpage
